%!TEX root =  main.tex


\para{Speculative decoding.} SD has been receiving increasing attention as an effective method to accelerate LLM inference by decomposing autoregressive decoding into a drafting phase and a verification phase~\citep{leviathan2023,ryu2024closerlook-survey,xia2024sd-survey}. 
Prior work has focused on improving SD performance in both single-request settings~\cite{cai2024medusa,cheng2024recurrentdrafterfastspeculative,zhang2024skiplayer,li2025eagle3,oliaro2025suffixdecoding} and batch-request settings~\citep{li2025adaserve,liu2025turbospec,wu2025tetris}, primarily by improving drafting efficiency. 
In contrast, \alg{} explores an orthogonal direction by parallelizing the drafting and verification phases, allowing easy integration with existing methods to further improve SD performance.


\para{Parallel speculative decoding.}
Recent work has explored reducing drafting overhead through parallelism. Parallel SD methods achieve faster decoding at the cost of substantial GPU resources, such as increased VRAM consumption~\citep{wang2024minions,timor2025distributed}. Another promising direction is parallel decoding with tree attention, which accelerates SD by drafting multiple tokens in parallel~\citep{cai2024medusa,li2024eagle}. 
In addition, prior work focuses on the single-request setting and modifies the SD process to generate variable-length drafts, thereby reducing synchronization overhead between the drafter and verifier~\citep{liu2025pearl}.
Whereas \alg{} adopts a new parallelism paradigm, which we term \textit{Batch Parallelism}. This paradigm requires only one additional GPU compared to standard SD, balances GPU resource usage and overhead reduction, and avoids the need for a dedicated training stage.


\para{Adaptive drafting.}
Under many concurrent requests, standard SD can have suboptimal performance when using long speculative lengths, as the increased number of draft tokens places significant pressure on the verification stage. To mitigate this issue, adaptive drafting strategies generate variable-length drafts and selectively verify only the most promising tokens. EAGLE-2~\citep{li2024eagle2} adopts this approach by expanding and reranking a context-aware draft tree. Other heuristics include selecting drafts with the highest likelihood of acceptance within a batch~\citep{wu2025tetris} and choosing among drafts generated by different speculative methods based on their historical acceptance rates~\citep{hou2025banditspec}. Additionally, SpecServe~\citep{huang2025specserve} dynamically controls the execution of the draft model based on estimated draft quality and generation speed.
Despite these advances, existing adaptive drafting methods yield limited performance gains: the number of draft tokens that can be generated remains constrained, as excessive drafting quickly makes verification overhead prohibitive.